  % ======================================================================
  %  Personal LaTeX Template
  %  Victor Gabriele De Matteis
  %
  %  Usage:
  %    1. Copy this file and rename it for your project.
  %    2. Fill in the title block below.
  %    3. Replace \bibliography{references} with your own .bib file.
  %    4. Compile: pdflatex -> bibtex -> pdflatex x2
  %
  %  Citation style: numeric [1], compact (abbrvnat).
  %  To switch to author-year, replace:
  %    \usepackage[numbers,square,sort&compress]{natbib}
  %    \bibliographystyle{abbrvnat}
  %  with:
  %    \usepackage[round,authoryear]{natbib}
  %    \bibliographystyle{apalike}
  % ======================================================================

  \documentclass[10pt,a4paper]{scrartcl}

  % --------------------------------------------------
  % Encoding and language
  % --------------------------------------------------
  \usepackage[T1]{fontenc}
  \usepackage[utf8]{inputenc}
  \usepackage[english]{babel}

  % --------------------------------------------------
  % Page layout
  % --------------------------------------------------
  \usepackage[
    top=18mm,
    bottom=18mm,
    left=20mm,
    right=20mm,
    includefoot
  ]{geometry}

  % --------------------------------------------------
  % Mathematics and graphics
  % --------------------------------------------------
  \usepackage{graphicx}
  \usepackage{amsmath,amssymb}
  \usepackage{bm}
  \setlength{\jot}{0.5cm}

  % --------------------------------------------------
  % Tables
  % --------------------------------------------------
  \usepackage{booktabs}
  \usepackage{tabularx}
  \usepackage{array}
  \usepackage{ragged2e}

  % A ragged-right fixed-width column for use in tabularx:
  %   \begin{tabularx}{\textwidth}{@{}L{3cm} X@{}}
  \newcolumntype{L}[1]{>{\RaggedRight\arraybackslash}p{#1}}

  % --------------------------------------------------
  % Typography
  % --------------------------------------------------
  \usepackage[scaled=0.95]{helvet}
  \usepackage[expansion=false]{microtype}   % expansion=false avoids font errors
  \renewcommand{\familydefault}{\sfdefault} % sans-serif throughout

  % --------------------------------------------------
  % Lists
  % --------------------------------------------------
  \usepackage{enumitem}
  \setlist[itemize]{
    leftmargin=*,
    topsep=2pt,
    parsep=0pt,
    partopsep=0pt,
    itemsep=3pt
  }
  \setlist[description]{
    leftmargin=0pt,
    topsep=2pt,
    parsep=0pt,
    itemsep=3pt
  }

  % --------------------------------------------------
  % Colour and hyperlinks
  % --------------------------------------------------
  \usepackage[dvipsnames]{xcolor}
  \usepackage[hidelinks]{hyperref}

  \definecolor{accent}{RGB}{0,129,88}   % change this to restyle everything

  % \emph produces upright accent-coloured text instead of italic.
  % Use \textit{} when you genuinely want italics.
  \renewcommand{\emph}[1]{\textcolor{accent}{#1}}

  % --------------------------------------------------
  % Citations
  % --------------------------------------------------
  \usepackage[numbers,square,sort&compress]{natbib}

  % --------------------------------------------------
  % Header and footer
  % --------------------------------------------------
  \usepackage[automark,headsepline=false,footsepline=false]{scrlayer-scrpage}
  \clearpairofpagestyles

  %  Edit the three footer fields below:
  \ifoot{Author Name}          % left:   your name / student ID
  \cfoot{\pagemark}            % centre: page number (always keep this)
  \ofoot{Course / Document}    % right:  course name or document label

  \setlength{\footskip}{20mm}  % pushes footer toward the page edge
  \pagestyle{scrheadings}

  % --------------------------------------------------
  % Spacing
  % --------------------------------------------------
  \setlength{\parindent}{0pt}
  \setlength{\parskip}{0.45em}

  % --------------------------------------------------
  % Section formatting
  % --------------------------------------------------
  \RedeclareSectionCommand[
    beforeskip=1.4ex plus 0.2ex minus 0.1ex,
    afterskip=0.6ex
  ]{section}
  \RedeclareSectionCommand[
    beforeskip=1.0ex plus 0.2ex minus 0.1ex,
    afterskip=0.4ex
  ]{subsection}
  \RedeclareSectionCommand[
    beforeskip=0.7ex plus 0.2ex minus 0.1ex,
    afterskip=0.3ex
  ]{subsubsection}

  \setkomafont{section}{\large\bfseries\color{accent}}
  \setkomafont{subsection}{\normalsize\bfseries\color{accent}}
  \setkomafont{subsubsection}{\normalsize\bfseries}

  % A section command that adds a faint rule beneath the title.
  % Use \RuleSection{Title} instead of \section{Title} where you want it.
  \newcommand{\RuleSection}[1]{%
    \section{#1}%
    \vspace{0.2em}{\color{accent!35}\hrule height 0.4pt}\vspace{0.3em}%
  }

  % ======================================================================
  \begin{document}
  % ======================================================================

  % --------------------------------------------------
  % Title block
  % --------------------------------------------------
  \begin{center}
    {\Huge\bfseries\color{accent} Document Title}\\[0.25cm]
    {\small Institution --- Course --- A.Y.\ 20XX/20XX}\\[0.15cm]
    {\small Author Name \quad\textbullet\quad Student ID}\\[0.2cm]
    {\color{accent!35}\hrule height 0.5pt}
  \end{center}

  \vspace{0.3em}

  % Optional abstract — delete if not needed.
  \begin{abstract}
  \noindent
  Abstract text goes here.
  \end{abstract}

  % ======================================================================
  %  BODY
  % ======================================================================

  \RuleSection{First Section}
  \label{sec:first}

  Body text. Use \emph{emphasis} for accent-coloured terms and \textit{italics}
  when you genuinely mean italics. Cite with \citep{key} for parenthetical
  references \citep{example} or \citet{key} for inline ones.

  \subsection{A Subsection}
  \label{sec:sub}

  More text. Cross-reference with Section~\ref{sec:first}; always use
  \verb|\label| and \verb|\ref|, never hardcode numbers.

  % --------------------------------------------------
  % Example: description list (feature list style)
  % --------------------------------------------------
  \begin{description}
    \item[] \textbf{Item A} --- description of item A.
    \item[] \textbf{Item B} --- description of item B.
  \end{description}

  % --------------------------------------------------
  % Example: table (use \\[.3cm] or more between rows)
  % --------------------------------------------------
  \begin{table}[h]
  \small
  \centering
  \renewcommand{\arraystretch}{1.1}
  \begin{tabularx}{\textwidth}{@{}L{3cm} X@{}}
  \toprule
  \textbf{Column A} & \textbf{Column B} \\[.3cm]
  \midrule
  Row 1 & Description of row 1. \\[.3cm]
  Row 2 & Description of row 2. \\[.3cm]
  \bottomrule
  \end{tabularx}
  \caption{A short, informative caption.}
  \label{tab:example}
  \end{table}

  % ======================================================================
  % Mathematical examples
  % ======================================================================

  \RuleSection{Mathematical Examples}

  % --------------------------------------------------
  \subsection{Aligned equations}

  A simple derivation can be written using the \texttt{align} environment:

  \begin{align}
      y &= mx + q \\
      &= m(x_1+x_2) + q \\
      &= mx_1 + mx_2 + q .
  \end{align}


  % --------------------------------------------------
  \subsection{Systems of equations}

  A linear system:

  \[
  \begin{cases}
      a_{11}x_1 + a_{12}x_2 = b_1,\\
      a_{21}x_1 + a_{22}x_2 = b_2 .
  \end{cases}
  \]


  A more compact matrix representation:

  \begin{equation}
      \bm{A}\bm{x}=\bm{b},
  \end{equation}

  where

  \[
  \bm{A}=
  \begin{bmatrix}
  a_{11} & a_{12}\\
  a_{21} & a_{22}
  \end{bmatrix},
  \qquad
  \bm{x}=
  \begin{bmatrix}
  x_1\\
  x_2
  \end{bmatrix}.
  \]


  % --------------------------------------------------
  \subsection{Long aligned derivation}

  For example, solving a quadratic equation:

  \begin{align}
      ax^2+bx+c &=0\\
      ax^2+bx &= -c\\
      x^2+\frac{b}{a}x &= -\frac{c}{a}\\
      x^2+\frac{b}{a}x+\frac{b^2}{4a^2}
          &= \frac{b^2}{4a^2}-\frac{c}{a}\\
      \left(x+\frac{b}{2a}\right)^2
          &=\frac{b^2-4ac}{4a^2}.
  \end{align}


  % --------------------------------------------------
  \subsection{Piecewise functions}

  \[
  f(x)=
  \begin{cases}
      x^2, & x\geq0,\\
      -x, & x<0 .
  \end{cases}
  \]


  % --------------------------------------------------
  \subsection{Optimization problem}

  A typical control/engineering formulation:

  \begin{equation}
      \begin{aligned}
      \min_{\bm{x}}\quad&
          J(\bm{x})
          =
          \sum_{k=0}^{N}
          \left\|
          \bm{x}_k-\bm{x}_{ref}
          \right\|^2_Q
          +
          \left\|
          \bm{u}_k
          \right\|^2_R
          \\
      \text{s.t.}\quad&
          \bm{x}_{k+1}
          =
          f(\bm{x}_k,\bm{u}_k),\\
      &
          \bm{u}_k\in\mathcal{U}.
      \end{aligned}
  \end{equation}


  % --------------------------------------------------
  \subsection{Control-system equations}

  State-space representation:

  \begin{align}
      \dot{\bm{x}}(t)
          &=A\bm{x}(t)+B\bm{u}(t),\\
      \bm{y}(t)
          &=C\bm{x}(t)+D\bm{u}(t).
  \end{align}


  % --------------------------------------------------
  \subsection{Probability and statistics}

  Expectation:

  \begin{equation}
      \mathbb{E}[X]
      =
      \int_{-\infty}^{+\infty}
      x f_X(x)\,dx .
  \end{equation}


  Variance:

  \begin{equation}
      \mathrm{Var}(X)
      =
      \mathbb{E}
      \left[
          (X-\mathbb{E}[X])^2
      \right].
  \end{equation}


  % --------------------------------------------------
  \subsection{Numbered theorem-like equation}

  \begin{equation}
      \|Ax\|
      \leq
      \|A\|\|x\|.
      \label{eq:norm_property}
  \end{equation}

  Equation~\eqref{eq:norm_property} is a standard norm inequality.

  % ======================================================================
  \bibliographystyle{abbrvnat}
  \bibliography{references}   % <- replace with your .bib filename (no extension)

  \end{document}